\subsection{自由模的线性映射及主理想整环上的矩阵}

设 A 为主理想整环。考虑从秩为 m 的自由 A-模 L 到秩为 n 的自由 A-模 \(L'\) 的一个线性映射 f。前面的结果使我们可以通过为 L 和 \(L'\) 选取适当的基，把 f 的矩阵化为一种特别简单的形式，称为矩阵的典范形式。

命题 5. --- 设 A 为主理想整环，f 为从秩为 m 的自由 A-模 L 到秩为 n 的自由 A-模 L' 的一个秩为 r 的线性映射。则存在 L 的基 \((e_{i})\) \((1 \leqslant i \leqslant m)\) 和 L' 的基 \((e_{j}')\) \((1 \leqslant j \leqslant n)\)，使得 \(f(e_{i}) = \alpha_{i} e_{i}'\)（对 \(1 \leqslant i \leqslant r\)），且 \(f(e_{i}) = 0\)（对 i \textgreater{} r），其中 \(\alpha_{i}\) 是 A 的非零元素，且每一个整除其后一个；诸理想 \(A\alpha_{i}\) 是 \(f(L)\) 在 \(L'\) 中的不变因子，因此是唯一确定的。

设 \(L_0 = f^{-1}(0)\) 是 \(f\) 的核；商 \(L/L_0\) 同构于模 \(f(L)\)，它是 \(L'\) 的子模，因而是自由的（VII，第 15 页，推论 2）；于是 \(L_0\) 有一个补 \(L_1\)（在 \(L\) 中；II，第 218 页，命题 21），且 \(f\) 在 \(L_1\) 上的限制是从 \(L_1\) 映到 \(f(L) = M'\) 的同构。如果诸理想 \(A\alpha_i\)（\(1 \leqslant i \leqslant r\)）是 \(M'\) 在 \(L'\) 中的不变因子，则 VII 第 18 页的定理 1 表明存在一组基 \((e_j')\)（\(1 \leqslant j \leqslant n\)）（\(L'\) 的基），使得 \((\alpha_i e_i')\)（\(1 \leqslant i \leqslant r\)）是 \(M'\) 的一组基。由于 \(f\) 在 \(L_1\) 上的限制是从 \(L_1\) 映到 \(M'\) 的同构，故存在一组基 \((e_i)\)（\(1 \leqslant i \leqslant r\)）（\(L_1\) 的基），使得 \(f(e_i) = \alpha_i e_i'\)。这组基可以扩充为 \((e_k)\)（\(1 \leqslant k \leqslant m\)）（即 \(L\) 的一组基），办法是取 \((e_s)\)（\(r + 1 \leqslant s \leqslant m\)）为核的基

\[
\mathbf {L} _ {0}
\]

推论 1. --- 设 X 是主理想整环 A 上的一个秩为 r 的矩阵，有 n 行和 m 列；则存在一个与 X 等价的矩阵 \(X_{0}\)（II，第 354 页），具有形式

\[
\left( \begin{array}{c c c c c c c} \alpha_ {1} & 0 & \ldots & 0 & 0 & \ldots & 0 \\ 0 & \alpha_ {2} & \ldots & 0 & 0 & \ldots & 0 \\ 0 & 0 & \ldots & \alpha_ {r} & 0 & \ldots & 0 \\ 0 & 0 & \ldots & 0 & 0 & \ldots & 0 \\ 0 & 0 & \ldots & 0 & 0 & \ldots & 0 \end{array} \right)
\]

其中 \(\alpha_{i}\) 是 A 的非零元素，且每个整除下一个。在这些条件下，诸 \(\alpha_{i}\) 在相差一个可逆元素乘法的意义下是唯一确定的。

鉴于两个矩阵 X 和 \(X'\) 等价是指存在 A 上的阶分别为 n 和 m 的可逆方阵 P 和 Q，使得 \(X' = PXQ\)，推论 1 正是用矩阵语言表述的命题 5。

在命题 5 和推论 1 的记号下，（非零）理想 \(\mathbf{A}\alpha_{i}\) 称为线性映射 \(f\) 或矩阵 \(X\) 的不变因子。于是由推论 1 立即得出：

推论 2. --- 设 X 和 \(X'\) 是主理想整环 A 上具有 n 行 m 列的两个矩阵，则它们等价当且仅当它们具有相同的不变因子。

注意，当 A 是域时，我们可以取 \(\alpha_{i}\) 等于 1，于是我们重新得到 II 第 360 页的命题 13。

若 X 是线性映射 f 关于 L 的任意一组基和 \(L'\) 的任意一组基的矩阵，则 X 的各列是 \(L'\) 中某些元素关于 \(L'\) 的基的坐标向量，而这些元素构成 \(f(\mathrm{L})\) 的一个生成元组。因此，下面的结果是命题 4 的直接推论。

命题 6. --- 设 X 是主理想整环 A 上秩为 r 的矩阵，并设 \(A\alpha_{i}\) ( \(1 \leqslant i \leqslant r\) ) 是其不变因子序列。则 \(\alpha_{1}\) 是 X 的元素的一个最大公因子；且乘积 \(\alpha_{1} \ldots \alpha_{q}\) 是 X 的 q 阶子式的一个最大公因子，其中 \(q \leqslant r\) .

